\section*{Chapitre \ref{ch:hf:options-cross-venue-and-volatility-arbitrage-at-speed} --- Options~: arbitrage inter-plateformes et de volatilité à la vitesse}

\begin{solution}{exo:hf:options-cross-venue-and-volatility-arbitrage-at-speed:1}
$-\ln(9.90/10)/0.25=4.02~\%$.
\end{solution}

\begin{solution}{exo:hf:options-cross-venue-and-volatility-arbitrage-at-speed:2}
$4.60-4.10-100.01+0.50+100e^{-0.01}=-0.005$~: un demi-cent négatif avant frais~; pas d'arbitrage.
\end{solution}

\begin{solution}{exo:hf:options-cross-venue-and-volatility-arbitrage-at-speed:3}
$c=4\times0.2+2\times0.1=1.0$ cent par contrat négocié (par action), $g=30$ cents~: le jeu rapporte au-delà de $c/(2g)=1/60\approx1.7~\%$.
\end{solution}

\begin{solution}{exo:hf:options-cross-venue-and-volatility-arbitrage-at-speed:4}
Les meilleurs cours acheteur et vendeur sur un plus grand nombre de plateformes sont plus extrêmes, donc ils franchissent la parité plus souvent~; mais chaque
violation est un bruit propre à une plateforme, que le rafraîchissement suivant efface, et tous les scanners voient les mêmes cotations~: seul le plus rapide la capture.
\end{solution}

\begin{solution}{exo:hf:options-cross-venue-and-volatility-arbitrage-at-speed:5}
Les bons du Trésor portent un rendement d'opportunité (convenience yield), lié à la liquidité, à la valeur de collatéral et au traitement réglementaire, qui abaisse leur rendement~; le taux de la box, déduit sans
ce terme, est le taux sans risque avant cette prime, soit environ 40 points de base de plus en moyenne chez van Binsbergen, Diamond et Grotteria.
\end{solution}

\begin{solution}{exo:hf:options-cross-venue-and-volatility-arbitrage-at-speed:6}
Il gonfle le volume le dernier jour de cotation avec dividende, grâce à des transactions qui ne transfèrent aucun risque~; quiconque lit le volume comme une liquidité ou un intérêt se trompe.
\end{solution}

\begin{solution}{exo:hf:options-cross-venue-and-volatility-arbitrage-at-speed:7}
Environ 13 cents.
\end{solution}

\begin{solution}{exo:hf:options-cross-venue-and-volatility-arbitrage-at-speed:8}
Les violations sont comptées avant exécution~: après un rafraîchissement, avec la moitié du bruit conservée, il n'en reste qu'environ 6~\%, et tous les autres scanners voient les mêmes
cotations. Le revenu est la part que la société gagne à la vitesse, nette de frais, et non le nombre de violations.
\end{solution}

\begin{solution}{pb:hf:options-cross-venue-and-volatility-arbitrage-at-speed:1}
\begin{enumerate}
\item Voir la \cref{def:hf:options-cross-venue-and-volatility-arbitrage-at-speed:conversion}.
\item Neuf strikes, plateformes cotant large de 10 cents avec 3 cents de bruit, un cent par jambe~: violations dans 0.7~\%, 6.6~\% et 26~\% des instantanés pour 1, 3 et 6
plateformes, d'une valeur d'un à deux cents.
\item 6 à 8~\% avec la moitié du bruit conservée, 43~\% avec 90~\%~: la persistance des cotations.
\item Le reversal est une vente à découvert de l'action qui paie les frais d'emprunt et les dividendes.
\item 4.02~\%~; un cent contre un taux de financement de 4~\%.
\item Rendements d'opportunité des bons du Trésor d'environ 40 points de base, plus élevés en dessous de trois mois, quatre fois plus élevés pendant la crise.
\item Voir la \cref{def:hf:options-cross-venue-and-volatility-arbitrage-at-speed:jelly}~; 4.0~\% de carry entre trois et six mois.
\item Une annonce de dividende ou une modification des frais d'emprunt entre les échéances.
\item Lorsque le dividende dépasse la valeur temps abandonnée (la valeur du put au strike)~; 13 cents dans l'exemple.
\item Plus de la moitié des positions longues qui auraient dû être exercées ne l'ont pas été, de 1996 à 2006, ce qui a coûté aux détenteurs plus de 491~millions de dollars.
\item Des transactions compensatoires entre teneurs de marché qui exercent toutes les positions longues et prennent la plupart des assignations~; le volume qu'elles gonflent, favorisé par les
plafonds de frais des bourses.
\item $c/(2g)$~: 1.7~\%.
\item Nul en dessous de 1.7~\%~; 0.27~\$ par contrat public à 5~\%, 1.05~\$ à 10~\%, 3.04~\$ à 20~\%, 10.02~\$ à 50~\%, capturant de 42~\% à 82~\% de la
valeur des défaillances.
\item Les détenteurs qui n'exercent pas, soit en exerçant eux-mêmes, soit par l'intermédiaire de leurs courtiers.
\item Des frais plus bas abaissent le seuil de rentabilité et augmentent la taille du jeu.
\item Les cotations d'une même série cotée sur un marché suivent celles de l'autre avec un retard~; la transaction prend la cotation périmée.
\item La box, qui immobilise du capital et du bilan pendant toute sa durée.
\item À partir des positions ouvertes sur les calls qui devraient être exercés, et des taux d'exercice passés par type de détenteur aux dates ex-dividende.
\item L'avis et, lorsqu'elles ont été convenues, les instructions d'exercice automatique.
\item Les détenteurs qui oublient, rémunérés par quiconque négocie pour reprendre leurs assignations.
\end{enumerate}
\end{solution}

\begin{solution}{iq:hf:options-cross-venue-and-volatility-arbitrage-at-speed:1}
$C-P=S-\mathrm{PV}(D)-Ke^{-rT}$~; si le call est cher, on vend le call, on achète le put et l'action (une conversion), si le dime dépasse les coûts.
\iqlookfor{parité avec dividendes, le bon trade.}
\end{solution}

\begin{solution}{iq:hf:options-cross-venue-and-volatility-arbitrage-at-speed:2}
On cote des box sur des options indicielles européennes à différentes échéances, on déduit les taux zéro, et on retranche les rendements des bons du Trésor de même échéance.
\iqlookfor{taux implicites des box.}
\end{solution}

\begin{solution}{iq:hf:options-cross-venue-and-volatility-arbitrage-at-speed:3}
La plupart seront exercées et on sera assigné~: on vend l'action à découvert le lendemain matin, sans le dividende. On achète l'action ou on exerce les positions longues compensatoires~;
la fraction non assignée gagne le dividende diminué de la valeur temps.
\iqlookfor{assignation et sa position.}
\end{solution}

\begin{solution}{iq:hf:options-cross-venue-and-volatility-arbitrage-at-speed:4}
On tient à jour, de façon incrémentale, le meilleur cours acheteur et le meilleur cours vendeur par série sur l'ensemble des plateformes~; on ne recalcule la parité que pour les séries dont les cotations ont changé~; on exige des quantités sur toutes les jambes
et on applique les frais, les frais d'emprunt et les dividendes~; on envoie toutes les jambes ensemble.
\iqlookfor{incrémental, exécutable, toutes les jambes.}
\end{solution}

\begin{solution}{iq:hf:options-cross-venue-and-volatility-arbitrage-at-speed:5}
Les frais d'emprunt grimpent ou les actions sont rappelées~: le reversal perd son carry ou doit être débouclé~; la parité de cette action se déplace et les puts deviennent chers.
\iqlookfor{risque d'emprunt.}
\end{solution}

\begin{solution}{iq:hf:options-cross-venue-and-volatility-arbitrage-at-speed:6}
La dérivée de $aq/(P+2q)$ vaut $aP/(P+2q)^2$~; en l'égalant à $c$, on obtient $(P+2q)^2=aP/c$, donc $q^\ast=(\sqrt{aP/c}-P)/2$.
\iqlookfor{la condition du premier ordre.}
\end{solution}
